% ============================================================
%  GUÍA DE PRÁCTICA — TEMA 1 (PRIMER LAPSO) — 3 SECCIONES
%  Naturales, Decimales, Operaciones Combinadas y Ecuaciones
%  Curso: 1er Año — CELAM
%  Autor: Ing. Gabriel Astudillo
%  Basada en examen_real_t1l1.md: cada sección repite la MISMA
%  estructura del examen (operaciones + combinada + ecuaciones)
%  y un nivel de dificultad equivalente.
%  Incluye clave de respuestas entre \ifmwclaves ... \fi
%  (genera versión normal y versión _scr sin respuestas).
% ============================================================

\documentclass[12pt, letterpaper]{article}

% ── Paquetes ─────────────────────────────────────────────────
\usepackage{mathwizards-guia}
\mwguia{Guía de Práctica — Tema 1 (3 Secciones)}{Ing. Gabriel Astudillo $\cdot$ 1er Año}

% ── Comandos propios de esta guía ────────────────────────────
\newcommand{\N}{\mathbb{N}}

\begin{document}
% ============================================================

% ── Portada / Título ─────────────────────────────────────────
\mwportada{Guía de Práctica — Tema 1}{Tres Secciones al Estilo del Examen}{Naturales, Decimales, Operaciones Combinadas y Ecuaciones}

\vspace{4pt}
\begin{cuadroinfo}
  \small
  \textbf{Presentación:} esta guía reúne \textbf{tres secciones} con la \emph{misma estructura del examen real del Tema 1}: en cada una resolverás (1)~una lista de operaciones con naturales y decimales, (2)~una operación combinada y (3)~ecuaciones de primer grado en una incógnita. Así practicas con \textbf{tres versiones} de dificultad equivalente antes de la evaluación. Al final encontrarás la \textbf{clave de respuestas}.
  \\[4pt]
  \textbf{Convenciones:}
  \begin{itemize}[leftmargin=*]
    \item Trabajamos en los naturales: si la solución de una ecuación \textbf{no pertenece a $\N$}, escribe $x \notin \N$ y no la compruebes.
    \item En las divisiones inexactas, calcula el cociente \textbf{hasta dos cifras decimales}.
    \item Usa la \emph{coma decimal} y cuida los \emph{ceros} no significativos.
    \item Muestra \emph{todos} los pasos: se evalúa el procedimiento.
  \end{itemize}
\end{cuadroinfo}

\vspace{8pt}

% ============================================================
\section{Sección 1}
% ============================================================

\subsection*{1)\quad Ordenar y resolver las siguientes operaciones. \hfill\normalfont\small (2{,}5 ptos c/u)}
\begin{enumerate}[label=\alph*), leftmargin=*]
  \item $425\,832 \div 24$
  \item $845 \cdot 726$
  \item $58\,742\,316 + 37\,649\,857$
  \item $83\,562\,471 - 27\,438\,596$
  \item $29\,040 \div 2{,}4$
  \item $38{,}5 \div 0{,}154$
\end{enumerate}

\vspace{6pt}

\subsection*{2)\quad Resolver la siguiente operación combinada. \hfill\normalfont\small (2 ptos)}
\begin{enumerate}[label=\alph*), leftmargin=*]
  \item $3 \cdot (5 + 3^2 \div 3^1) + 4^2 \cdot (9 - 2^3) + 2^3 \cdot (36 \div 4 - 45 \div 9)$
\end{enumerate}

\vspace{6pt}

\subsection*{3)\quad Hallar el valor de la incógnita en las siguientes ecuaciones. \hfill\normalfont\small (3 ptos)}
\begin{enumerate}[label=\alph*), leftmargin=*]
  \item $\dfrac{3 \cdot (2x - 1)}{4} = \dfrac{2 \cdot (x + 3)}{3}$
  \item $\dfrac{2 \cdot (3x + 1)}{5} = \dfrac{4 \cdot (x - 1)}{3}$
\end{enumerate}

\newpage

% ============================================================
\section{Sección 2}
% ============================================================

\subsection*{1)\quad Ordenar y resolver las siguientes operaciones. \hfill\normalfont\small (2{,}5 ptos c/u)}
\begin{enumerate}[label=\alph*), leftmargin=*]
  \item $312\,768 \div 32$
  \item $654 \cdot 879$
  \item $63\,574\,218 + 28\,736\,459$
  \item $91\,235\,647 - 46\,789\,352$
  \item $58\,500 \div 7{,}5$
  \item $12{,}6 \div 0{,}36$
\end{enumerate}

\vspace{6pt}

\subsection*{2)\quad Resolver la siguiente operación combinada. \hfill\normalfont\small (2 ptos)}
\begin{enumerate}[label=\alph*), leftmargin=*]
  \item $4 \cdot (6 + 5^2 \div 5^1) + 3^2 \cdot (12 - 2^3) + 5^2 \cdot (48 \div 6 - 40 \div 8)$
\end{enumerate}

\vspace{6pt}

\subsection*{3)\quad Hallar el valor de la incógnita en las siguientes ecuaciones. \hfill\normalfont\small (3 ptos)}
\begin{enumerate}[label=\alph*), leftmargin=*]
  \item $\dfrac{3 \cdot (2x - 1)}{5} = \dfrac{2 \cdot (x + 3)}{4}$
  \item $\dfrac{5 \cdot (2x - 1)}{4} = \dfrac{3 \cdot (x + 2)}{2}$
\end{enumerate}

\newpage

% ============================================================
\section{Sección 3}
% ============================================================

\subsection*{1)\quad Ordenar y resolver las siguientes operaciones. \hfill\normalfont\small (2{,}5 ptos c/u)}
\begin{enumerate}[label=\alph*), leftmargin=*]
  \item $546\,336 \div 48$
  \item $937 \cdot 826$
  \item $69\,418\,523 + 27\,365\,489$
  \item $72\,543\,816 - 38\,697\,425$
  \item $46\,800 \div 5{,}2$
  \item $17{,}5 \div 0{,}35$
\end{enumerate}

\vspace{6pt}

\subsection*{2)\quad Resolver la siguiente operación combinada. \hfill\normalfont\small (2 ptos)}
\begin{enumerate}[label=\alph*), leftmargin=*]
  \item $2^2 \cdot (7 + 6^2 \div 3^2) + 3^1 \cdot (15 - 2^3) + 4^2 \cdot (54 \div 6 - 48 \div 8)$
\end{enumerate}

\vspace{6pt}

\subsection*{3)\quad Hallar el valor de la incógnita en las siguientes ecuaciones. \hfill\normalfont\small (3 ptos)}
\begin{enumerate}[label=\alph*), leftmargin=*]
  \item $\dfrac{4 \cdot (2x - 1)}{5} = \dfrac{3 \cdot (x + 8)}{4}$
  \item $\dfrac{3 \cdot (2x - 1)}{4} = \dfrac{5 \cdot (x - 1)}{2}$
\end{enumerate}

% ============================================================
\ifmwclaves
  \newpage
  \section{Clave de Respuestas}
  % ============================================================

  \subsection*{Sección 1}
  \begin{enumerate}[leftmargin=*, label=\textbf{\arabic*.}]
    \item
          $a)\; 425\,832 \div 24 = 17\,743$; \quad
          $b)\; 845 \cdot 726 = 613\,470$; \quad
          $c)\; 58\,742\,316 + 37\,649\,857 = 96\,392\,173$; \\
          $d)\; 83\,562\,471 - 27\,438\,596 = 56\,123\,875$; \quad
          $e)\; 29\,040 \div 2{,}4 = 12\,100$; \quad
          $f)\; 38{,}5 \div 0{,}154 = 250$.
    \item $3 \cdot (5 + 9 \div 3) + 16 \cdot (9 - 8) + 8 \cdot (9 - 5) = 3 \cdot 8 + 16 \cdot 1 + 8 \cdot 4 = 24 + 16 + 32 = 72$.
    \item
          $a)\; x = \dfrac{33}{10} \notin \N$; \quad
          $b)\; x = 13$ (se comprueba: $2(3 \cdot 13 + 1)/5 = 4(13 - 1)/3 = 16$).
  \end{enumerate}

  \subsection*{Sección 2}
  \begin{enumerate}[leftmargin=*, label=\textbf{\arabic*.}]
    \item
          $a)\; 312\,768 \div 32 = 9\,774$; \quad
          $b)\; 654 \cdot 879 = 574\,866$; \quad
          $c)\; 63\,574\,218 + 28\,736\,459 = 92\,310\,677$; \\
          $d)\; 91\,235\,647 - 46\,789\,352 = 44\,446\,295$; \quad
          $e)\; 58\,500 \div 7{,}5 = 7\,800$; \quad
          $f)\; 12{,}6 \div 0{,}36 = 35$.
    \item $4 \cdot (6 + 25 \div 5) + 9 \cdot (12 - 8) + 25 \cdot (8 - 5) = 4 \cdot 11 + 9 \cdot 4 + 25 \cdot 3 = 44 + 36 + 75 = 155$.
    \item
          $a)\; x = 3$ (se comprueba: $3(2 \cdot 3 - 1)/5 = 2(3 + 3)/4 = 3$); \quad
          $b)\; x = \dfrac{17}{4} \notin \N$.
  \end{enumerate}

  \subsection*{Sección 3}
  \begin{enumerate}[leftmargin=*, label=\textbf{\arabic*.}]
    \item
          $a)\; 546\,336 \div 48 = 11\,382$; \quad
          $b)\; 937 \cdot 826 = 773\,962$; \quad
          $c)\; 69\,418\,523 + 27\,365\,489 = 96\,784\,012$; \\
          $d)\; 72\,543\,816 - 38\,697\,425 = 33\,846\,391$; \quad
          $e)\; 46\,800 \div 5{,}2 = 9\,000$; \quad
          $f)\; 17{,}5 \div 0{,}35 = 50$.
    \item $4 \cdot (7 + 36 \div 9) + 3 \cdot (15 - 8) + 16 \cdot (9 - 6) = 4 \cdot 11 + 3 \cdot 7 + 16 \cdot 3 = 44 + 21 + 48 = 113$.
    \item
          $a)\; x = 8$ (se comprueba: $4(2 \cdot 8 - 1)/5 = 3(8 + 8)/4 = 12$); \quad
          $b)\; x = \dfrac{7}{4} \notin \N$.
  \end{enumerate}

  \vspace{6pt}
  \begin{cuadroinfo}
    \small
    \textbf{Recuerda:} en las ecuaciones que \emph{sí} tienen solución natural, sustituye el valor para comprobar la igualdad; en las que no, basta con escribir $x \notin \N$.
  \end{cuadroinfo}
\fi
% ============================================================

\vspace{10pt}
\begin{cuadroinfo}
  \small
  \textbf{¡Mucho éxito!} Repite cada sección cronometrándote: son tres oportunidades para llegar al examen con la mecánica de las operaciones bien afinada.
\end{cuadroinfo}

\end{document}
